% Unit under test. Identification data comes from metadata.tex.

\FunctionalDescription{
  \TheUnit\ stores the heater temperature setpoint, checks it against the
  allowed range 0--100 and restores the last valid setpoint from NvM at
  start-up. Details: \DocRef{design}.
}

% {name}{direction}{type / range}{description}
\Interface{SetpointCtrl_Init()}{in}{void}{Initialises the unit, reads NvM}
\Interface{SetpointCtrl_Set(sp)}{in}{uint8, [0, 100]}{Sets a new setpoint, returns E_OK / E_NOT_OK}
\Interface{SetpointCtrl_GetState()}{out}{enum}{IDLE / DEGRADED}

% [component]{dependency}{replaced by}{behaviour in test}
% [component] is optional: functions of one component share a box in the diagrams.
\Dependency[NvM]{NvM_ReadBlock()}{Mock}{Returns a configurable value, records calls}
\Dependency[NvM]{NvM_WriteBlock()}{Mock}{Records written data}
\Dependency[Dem]{Dem_ReportError()}{Mock}{Records reported events}

% Who uses the unit: {component}{interfaces it calls}
% Callers + dependencies give the context diagram; dependencies + interfaces
% give the test harness diagram in the Test Environment chapter.
\Caller{HeaterApp}{SetpointCtrl_Init(), SetpointCtrl_Set()}
\Caller{DiagService}{SetpointCtrl_GetState()}

% Free text: any LaTeX, printed at the end of this chapter (Unit Under Test).
% Write \end{FreeText} on its own line.
\begin{FreeText}
\subsection{Public API (excerpt)}
From \texttt{include/setpoint\_ctrl.h} -- 50\% of the API shown:
\begin{verbatim}
Std_ReturnType SetpointCtrl_Init(void);
Std_ReturnType SetpointCtrl_Set(uint8 sp);   /* 0..100 */
SetpointCtrl_StateType SetpointCtrl_GetState(void);
\end{verbatim}
\end{FreeText}

% Free text example 2: own subsections, numbered formulas, cross-references.
% Normal LaTeX rules apply inside FreeText (write \% for a percent sign).
\begin{FreeText}
\subsection{Setpoint Processing}
Summary of the processing the test cases check against; see \DocRef{design}
for the full design.

\subsubsection{Range Limitation}
A requested setpoint $sp_{\mathrm{req}}$ is accepted only inside the valid range:
\begin{equation}
  sp_k =
  \begin{cases}
    sp_{\mathrm{req}} & \text{if } 0 \le sp_{\mathrm{req}} \le 100,\\
    sp_{k-1}          & \text{otherwise (request rejected, \texttt{E\_NOT\_OK})}.
  \end{cases}
  \label{eq:range}
\end{equation}

\subsubsection{Ramp Limitation}
The effective setpoint follows the accepted one with at most
$\Delta_{\max} = 2\,\mathrm{K/s}$, sampled every $T_s = 10\,\mathrm{ms}$:
\begin{align}
  \Delta_k     &= \operatorname{clamp}\!\left(sp_k - sp^{\mathrm{eff}}_{k-1},\; -\Delta_{\max} T_s,\; \Delta_{\max} T_s\right), \label{eq:ramp-step}\\
  sp^{\mathrm{eff}}_k &= sp^{\mathrm{eff}}_{k-1} + \Delta_k . \label{eq:ramp}
\end{align}
A step from 20 to 100 therefore takes $\lceil 80 / (\Delta_{\max} T_s) \rceil = 4000$
cycles ($40\,\mathrm{s}$). Equation \eqref{eq:range} defines the expected results of
TC-SPC-002; the ramp \eqref{eq:ramp-step}--\eqref{eq:ramp} (SWR-0130) is verified
in integration test.
\end{FreeText}
